\lecture{Лекция 6}{}

\subsection{Разновидности методов второго порядка}
\label{subsec:newton-variants}

Общая схема
\[
  x_{k+1} = x_k + \alpha_k p_k ,
  \qquad\text{где}\qquad
  H(x_k)\, p_k = -\nabla f(x_k)
\]
(ньютоновское направление). Методы различаются выбором шага $\alpha_k$.

\begin{enumerate}[label=\arabic*),left=0pt]
  \item $\alpha_k \equiv 1$ — \textbf{чистый метод Ньютона}.
  \item $\alpha_0 = 1$, $\lambda \in (0,1)$, $\eps \in \big(0, \tfrac12\big]$ —
        принцип дробления по условию
        \begin{equation}\label{eq:newton-armijo}
          f(x_k + \alpha p_k) - f(x_k)
          \ \le\ \alpha\,\eps\,\nabla\T f(x_k)\, p_k
        \end{equation}
        — \textbf{демпфированный метод Ньютона}.
  \item $\big(f(x_k + \alpha p_k) - f(x_k)\big) \to \min\limits_{\alpha > 0}
        \iff f(x_k + \alpha p_k) \to \min\limits_{\alpha > 0}$ —
        наискорейший спуск вдоль ньютоновского направления.
\end{enumerate}

\inputfig{figures/fig-newton-quadratic}

\begin{idea}[title={Почему единичный шаг}]
Ньютоновское направление $p_k$ ведёт ровно в вершину квадратичной
аппроксимации $\psi_k$, построенной в точке $x_k$:
\[
  x_k + 1 \cdot p_k = \argmin_x \psi_k(x) .
\]
То есть $\alpha_k = 1$ — это «пройти весь путь до минимума параболоида».
Вблизи $x_*$ аппроксимация почти точна, и единичный шаг даёт квадратичную
сходимость. Вдали от минимума $\psi_k$ может плохо описывать $f$, и
единичный шаг способен увеличить $f$ — отсюда вариант~2) с дроблением.

Именно это и показывает рисунок с вложенными параболами: на каждой итерации
строится своя квадратичная аппроксимация, её вершина становится следующей
точкой, и чем ближе к минимуму, тем теснее парабола прилегает к $f$.
\end{idea}

\begin{idea}[title={Почему $\eps \le \tfrac12$}]
Условие~\eqref{eq:newton-armijo} — обычное условие достаточного убывания, но
записанное через $\nabla\T f(x_k) p_k < 0$ (а не через $\norm{\nabla f}^2$,
как в градиентных методах): теперь направление не антиградиентное, и убывание
измеряется вдоль него.

Ограничение $\eps \le \tfrac12$ нужно, чтобы единичный шаг проходил проверку.
Посчитаем для квадратичной $f$. Так как $H p = -\nabla f$, имеем
$p\T H p = -\nabla\T f\, p$, поэтому
\[
  f(x_k + \alpha p) - f(x_k)
  = \alpha\,\nabla\T f\, p + \frac{\alpha^2}{2}\, p\T H p
  = \nabla\T f\, p \left(\alpha - \frac{\alpha^2}{2}\right) .
\]
При $\alpha = 1$ это $\tfrac12 \nabla\T f\, p$, а правая
часть~\eqref{eq:newton-armijo} равна $\eps\,\nabla\T f\, p$. Поскольку
$\nabla\T f\, p < 0$, условие
\[
  \tfrac12 \nabla\T f\, p \ \le\ \eps\,\nabla\T f\, p
  \qquad\Longleftrightarrow\qquad
  \eps \ \le\ \tfrac12 .
\]
То есть при $\eps > \tfrac12$ дробление отвергало бы ньютоновский шаг даже
там, где он идеален, и метод терял бы квадратичную скорость, скатываясь к
линейной.
\end{idea}

\subsubsection*{Модификации: экономия на гессиане}

\begin{idea}[title={Откуда этот материал}]
На лекции перечислены только три способа выбора шага. Ниже — две стандартные
модификации, снимающие главный недостаток метода: пересчёт и обращение $H(x_k)$
на каждой итерации. Третья модификация — регуляризация
$H_k + \lambda_k E$ — разбирается отдельно как метод Марквардта
(\S\ref{subsec:marquardt}).
\end{idea}

\paragraph{Метод с фиксированной матрицей.}
Гессиан вычисляется один раз, в начальной точке:
\[
  H_0 = \nabla^2 f(x_0) ,
  \qquad
  x_{k+1} = x_k - \alpha_k H_0^{-1}\, \nabla f(x_k) .
\]
Система $H_0 p = -\nabla f(x_k)$ решается с одним и тем же разложением
матрицы, поэтому каждая итерация стоит как у градиентного метода.

\paragraph{Метод с периодическим обновлением матрицы.}
Гессиан фиксируется на $t$ итераций подряд, затем пересчитывается:
\[
  x_{lt+i+1} = x_{lt+i} - \alpha_{lt+i}\, H^{-1}(x_{lt})\, \nabla f(x_{lt+i}) ,
  \qquad i = \overline{0,\,t-1} ,
\]
где $l$ — номер «блока» итераций. При $t = 1$ это обычный метод Ньютона, при
$t = \infty$ — метод с фиксированной матрицей.

\begin{idea}
Чем платим. Обе схемы — это метод вида $p_k = -F \nabla f(x_k)$ из
\S\ref{subsec:descent-direction} с матрицей $F$, которая \emph{не} равна
$H^{-1}(x_k)$. Направление остаётся направлением спуска, пока $F \succ 0$ (а
$H_0^{-1} \succ 0$ в силу сильной выпуклости), поэтому сходимость не теряется.
Теряется квадратичная скорость: оценка теоремы~\ref{thm:newton-rate} опиралась
на то, что квадратичная модель построена \emph{в текущей} точке. С устаревшей
матрицей скорость становится линейной, причём тем хуже, чем дальше $x_k$ ушла
от $x_{lt}$.

Разумный компромисс — умеренное $t$: стоимость итерации падает в разы, а
скорость остаётся сверхлинейной на участке, где кривизна меняется слабо. Та же
идея, доведённая до конца, приводит к квазиньютоновским методам
(\S\ref{subsec:quasi-newton-anchor}): там матрица не замораживается, а дёшево
подправляется по наблюдаемым парам «сдвиг аргумента — сдвиг градиента».
\end{idea}

\subsection{Сходимость метода второго порядка}

\begin{theorem}\label{thm:newton-rate}
Пусть $f(x)$ дважды непрерывно дифференцируема и удовлетворяет условиям:
\begin{enumerate}[label=\arabic*),left=0pt]
  \item $\exists\, m, M$ $(0 < m \le M)$:
        $\ m\norm{z}^2 \le z\T H(y)\, z \le M\norm{z}^2$ \ $\forall z, y$;
  \item $\exists\, R > 0$:
        $\ \norm{H(x) - H(y)} \le R\,\norm{x - y}$ \ $\forall x, y$.
\end{enumerate}
Пусть $\{x_k\}$ строится методом второго порядка с выбором шага по методу
дробления. Тогда
\[
  \{x_k\} \xrightarrow[k \to \infty]{} \argmin f(x)
\]
с \textbf{квадратичной} скоростью независимо от выбора начального приближения
$x_0$, причём
\[
  \norm{x_{k+1} - x_*} \ \le\ \frac{R}{m}\,\norm{x_k - x_*}^2 .
\]
\end{theorem}

\begin{idea}
Условие~1) — то же, что в теореме~\ref{thm:grad-rate}: оно даёт существование и
единственность $x_*$, а заодно обратимость $H$ (значит ньютоновское направление
определено и является направлением спуска). Условие~2) — новое: липшицевость
\emph{гессиана}. Именно она отвечает за квадратичную оценку: ошибка квадратичной
аппроксимации имеет порядок $R\norm{x - x_*}^2$, и ровно столько же остаётся от
расстояния до минимума после шага.

Глобальность сходимости обеспечивает дробление: вдали от $x_*$ шаг уменьшается
и $f$ монотонно убывает, а вблизи — по предыдущему пояснению единичный шаг
проходит проверку при $\eps \le \tfrac12$, и включается квадратичная скорость.
Это и есть главный выигрыш по сравнению с методами первого порядка, где
оценки~\eqref{eq:q-min} и теоремы~\ref{thm:const-step} давали лишь линейную
скорость с $q$, зависящим от обусловленности $m/M$.
\end{idea}
